% Nguồn SGK Toán 9 Tập một — Bài tập cuối chương II, trang 42--43:
% https://cdn3.olm.vn/upload/thuvienso/4491668113-page-42-1773022934501.png
% https://cdn3.olm.vn/upload/thuvienso/4491668113-page-43-1773022934822.png
%
% Nguồn SBT Toán 9 KNTT Tập một — Ôn tập chương II, trang 28--29:
% SBT TOAN 9 KNTT TAP 1.pdf
% SHA-256: d7ffd04618456682795977e5c47de7874161a35e8606e5bac1906881603a72f7

\section*{Bài tập cuối chương II}

\begin{exercises}
    \subsection{Bài tập}

    \textbf{A. TRẮC NGHIỆM}

    \begin{questions}[resume=SGK]
        \item Nghiệm của bất phương trình $-2x+1<0$ là
        \begin{tasks}[label=\Alph*.](4)
            \task $x<\dfrac{1}{2}$.
            \task $x>\dfrac{1}{2}$.
            \task $x\leq\dfrac{1}{2}$.
            \task $x\geq\dfrac{1}{2}$.
        \end{tasks}

        \item Điều kiện xác định của phương trình $\dfrac{x}{2x+1}+\dfrac{3}{x-5}=\dfrac{x}{(2x+1)(x-5)}$ là
        \begin{tasks}[label=\Alph*.](2)
            \task $x\ne-\dfrac{1}{2}$.
            \task $x\ne-\dfrac{1}{2}$ và $x\ne-5$.
            \task $x\ne5$.
            \task $x\ne-\dfrac{1}{2}$ và $x\ne5$.
        \end{tasks}

        \item Phương trình $x-1=m+4$ có nghiệm lớn hơn $1$ với
        \begin{tasks}[label=\Alph*.](4)
            \task $m\geq-4$.
            \task $m\leq4$.
            \task $m>-4$.
            \task $m<-4$.
        \end{tasks}

        \item Nghiệm của bất phương trình $1-2x\geq2-x$ là
        \begin{tasks}[label=\Alph*.](4)
            \task $x>\dfrac{1}{2}$.
            \task $x<\dfrac{1}{2}$.
            \task $x\leq-1$.
            \task $x\geq-1$.
        \end{tasks}

        \item Cho $a>b$. Khi đó ta có:
        \begin{tasks}[label=\Alph*.](2)
            \task $2a>3b$.
            \task $2a>2b+1$.
            \task $5a+1>5b+1$.
            \task $-3a<-3b-3$.
        \end{tasks}
    \end{questions}

    \textbf{B. TỰ LUẬN}

    \begin{questions}[resume=SGK]
        \item Giải các phương trình sau:
        \begin{questions}
            \item $(3x-1)^2-(x+2)^2=0$;

            \answerline
            \answerline

            \item $x(x+1)=2(x^2-1)$.

            \answerline
            \answerline
        \end{questions}

        \item Giải các phương trình sau:
        \begin{questions}
            \item $\dfrac{x}{x-5}-\dfrac{2}{x+5}=\dfrac{x^2}{x^2-25}$;

            \answerline
            \answerline
            \answerline

            \item $\dfrac{1}{x+1}-\dfrac{x}{x^2-x+1}=\dfrac{3}{x^3+1}$.

            \answerline
            \answerline
            \answerline
        \end{questions}

        \item Cho $a<b$, hãy so sánh:
        \begin{questions}
            \item $a+b+5$ với $2b+5$;

            \answerline
            \answerline

            \item $-2a-3$ với $-(a+b)-3$.

            \answerline
            \answerline
        \end{questions}

        \item Giải các bất phương trình:
        \begin{questions}
            \item $2x+3(x+1)>5x-(2x-4)$;

            \answerline
            \answerline

            \item $(x+1)(2x-1)<2x^2-4x+1$.

            \answerline
            \answerline
        \end{questions}

        \item Một hãng viễn thông nước ngoài có hai gói cước như sau:

        \begin{center}
        \begin{tblr}{
            width=\linewidth,
            colspec={X[l] X[l]},
            rows={valign=m},
            row{1}={font=\bfseries, halign=c}
        }
            Gói cước A & Gói cước B \\
            Cước thuê bao hằng tháng 32 USD & Cước thuê bao hằng tháng là 44 USD \\
            45 phút miễn phí & Không có phút miễn phí \\
            0,4 USD cho mỗi phút thêm & 0,25 USD/phút \\
        \end{tblr}
        \end{center}

        \begin{questions}
            \item Hãy viết một phương trình xác định thời gian gọi (phút) mà phí phải trả trong cùng một tháng của hai gói cước là như nhau và giải phương trình đó.

            \answerline
            \answerline

            \item Nếu khách hàng chỉ gọi tối đa là 180 phút trong 1 tháng thì nên dùng gói cước nào? Nếu khách hàng gọi 500 phút trong 1 tháng thì nên dùng gói cước nào?

            \answerline
            \answerline
        \end{questions}

        \item Thanh tham dự một kì kiểm tra năng lực tiếng Anh gồm 4 bài kiểm tra nghe, nói, đọc và viết. Mỗi bài kiểm tra có điểm là số nguyên từ 0 đến 10. Điểm trung bình của ba bài kiểm tra nghe, nói, đọc của Thanh là 6,7. Hỏi bài kiểm tra viết của Thanh cần được bao nhiêu điểm để điểm trung bình cả 4 bài kiểm tra được từ 7,0 trở lên? Biết điểm trung bình được tính gần đúng đến chữ số thập phân thứ nhất.

        \answerline
        \answerline
        \answerline

        \item Để lập đội tuyển năng khiếu về bóng rổ của trường, thầy thể dục đưa ra quy định tuyển chọn như sau: mỗi bạn dự tuyển sẽ được ném 15 quả bóng vào rổ, quả bóng vào rổ được cộng 2 điểm; quả bóng ném ra ngoài bị trừ 1 điểm. Nếu bạn nào có số điểm từ 15 điểm trở lên thì sẽ được chọn vào đội tuyển. Hỏi một học sinh muốn được chọn vào đội tuyển thì phải ném ít nhất bao nhiêu quả vào rổ?

        \answerline
        \answerline
        \answerline
    \end{questions}
\end{exercises}

\section*{Ôn tập chương II}

\begin{practice}
    \subsection{Luyện tập}

    \textbf{A. CÂU HỎI (Trắc nghiệm)}

    Chọn một phương án đúng trong mỗi câu sau:

    \textbf{Câu 1.} Nghiệm của bất phương trình $-5x-1<0$ là
    \begin{tasks}[label=\Alph*.](4)
        \task $x>-\dfrac{1}{5}$.
        \task $x<-\dfrac{1}{5}$.
        \task $x\geq-\dfrac{1}{5}$.
        \task $x\leq-\dfrac{1}{5}$.
    \end{tasks}

    \textbf{Câu 2.} Điều kiện xác định của phương trình $\dfrac{2x}{x-1}+\dfrac{3x+1}{2x+1}=\dfrac{7x^2}{(x-1)(2x+1)}$ là
    \begin{tasks}[label=\Alph*.](2)
        \task $x\ne1$.
        \task $x\ne-\dfrac{1}{2}$.
        \task $x\ne1$ và $x\ne-\dfrac{1}{2}$.
        \task $x\in\mathbb{R}$.
    \end{tasks}

    \textbf{Câu 3.} Phương trình $2x+1=m$ có nghiệm lớn hơn $-2$ với
    \begin{tasks}[label=\Alph*.](4)
        \task $m>0$.
        \task $m>-2$.
        \task $m>-3$.
        \task $m\leq-3$.
    \end{tasks}

    \textbf{Câu 4.} Nghiệm của bất phương trình $3x-1\leq2x+2$ là
    \begin{tasks}[label=\Alph*.](4)
        \task $x>3$.
        \task $x<3$.
        \task $x\geq3$.
        \task $x\leq3$.
    \end{tasks}

    \textbf{Câu 5.} Cho $a>b$, khi đó ta có
    \begin{tasks}[label=\Alph*.](2)
        \task $2a>b+1$.
        \task $-2a>-2b$.
        \task $2a>a+b$.
        \task $3a<a+2b$.
    \end{tasks}

    \textbf{B. BÀI TẬP}

    \begin{questions}[resume=SBT]
        \item Giải các phương trình sau:
        \begin{questions}
            \item $5x(x+2)-10x-20=0$;

            \answerline
            \answerline

            \item $x^2-4x=x-4$.

            \answerline
            \answerline
        \end{questions}

        \item Giải các phương trình sau:
        \begin{questions}
            \item $\dfrac{6}{8+x^3}=\dfrac{1}{x^2-2x+4}+\dfrac{1}{x+2}$;

            \answerline
            \answerline
            \answerline

            \item $\dfrac{x}{x+5}+\dfrac{x-5}{x}=2$.

            \answerline
            \answerline
            \answerline
        \end{questions}

        \item Giải các bất phương trình sau:
        \begin{questions}
            \item $(3x+1)(x+2)>x(3x-2)+1$;

            \answerline
            \answerline

            \item $2x(x+1)+3<x(2x+5)-7$.

            \answerline
            \answerline
        \end{questions}

        \item
        \begin{questions}
            \item Cho $a<b$ và $c<d$, chứng minh rằng $a+c<b+d$.

            \answerline
            \answerline
            \answerline

            \item Cho $0<a<b$ và $0<c<d$, chứng minh rằng $0<ac<bd$.

            \answerline
            \answerline
            \answerline
        \end{questions}

        \item Chứng minh rằng với số $a>0$, $b>0$ bất kì, ta luôn có $\dfrac{a}{b}+\dfrac{b}{a}\geq2$.

        \answerline
        \answerline
        \answerline

        \item Hà phải làm 4 bài kiểm tra tiếng Anh: nghe, nói, đọc và viết. Bài nghe, nói, đọc Hà đạt điểm số lần lượt là 78, 83 và 89. Hỏi bài kiểm tra viết, Hà phải đạt điểm số là bao nhiêu để điểm số trung bình Hà đạt được của cả 4 bài kiểm tra ít nhất là 85?

        \answerline
        \answerline
        \answerline

        \item Mức lương tối thiểu theo quy định ở Pháp năm 2022 là 10,25 € cho mỗi giờ làm việc. Trong dịp hè, Laurent David làm thêm tại một khách sạn theo mức lương tối thiểu như quy định và anh ấy muốn kiếm được ít nhất 1500 € trong mùa hè này.
        \begin{questions}
            \item Hãy viết một bất phương trình mô tả tình huống này.

            \answerline
            \answerline

            \item Hỏi anh ấy cần làm việc ít nhất bao nhiêu giờ để kiếm được số tiền trên?

            \answerline
            \answerline
        \end{questions}

        (€ là viết tắt của Euro, là loại tiền tệ mà 20 nước thuộc Liên minh Châu Âu đang sử dụng chung)

        \item Hai ô tô khởi hành cùng một lúc từ A đến B trên cùng quãng đường dài 150 km. Vận tốc xe thứ nhất hơn vận tốc xe thứ hai là 10 km/h và xe thứ nhất đến B sớm hơn xe thứ hai là 30 phút. Hỏi vận tốc của hai xe là bao nhiêu?

        \answerline
        \answerline
        \answerline
    \end{questions}
\end{practice}
